\chapter{The scene runtime}\label{ch:runtime}

The scene runtime executes a scene. It takes the validated record of Chapter~\ref{ch:erisml}, keeps
a small set of state machines and stratifications, and steps them on a stream of events. After every
event it reports which of the agent's actions are prohibited, which are obliged and which are
allowed. This chapter follows one step through the runtime, then describes the machines, the rule
by which prohibitions win, the discharge of obligations, the loading of strata, the scene agent whose
two models read and write the runtime, and the processes that present a scene's protocols as BPMN
diagrams. All of it is in the erisml-compiler package, under \code{runtime/}, \code{fsm/} and
\code{process/}. The twin of Part~III uses it unchanged.
% src: erisml-compiler src/erisml_compiler/runtime/scene.py:1-37
% src: gtc-prototype twin/brain.py:191-194, 312-320

\section{A runtime is built from one record}\label{sec:runtime-build}

The class \code{SceneRuntime} is constructed from a \code{CompilerIR} and an optional agent name.
When no name is passed it takes the scene's \code{extra.agent}. The constructor creates one
commitment machine per commitment. It creates one consent machine per stakeholder whose consent
status is set and is not \code{n/a}, starting in the state that status names, or in
\code{not\_obtained} for any other value. It creates one legitimacy machine per stakeholder in the
authority or protector role. It loads and checks the stratifications
(Section~\ref{sec:runtime-strata}). It then steps the scene's standing events in order of time
index.
% src: erisml-compiler src/erisml_compiler/runtime/scene.py:48, 106-135

For Margaret's home this gives three commitment machines, one consent machine for Margaret, three
legitimacy machines for the robot, the monitoring centre and emergency services, and seven
stratifications, fourteen state entries in all. None of the four standing events moves any of them.
The runtime therefore starts with every commitment active, Margaret's consent not obtained, every
authority fully legitimate and every stratification in its initial state.
% src: gtc-prototype twin/scene/margaret_home.erisml:41-69 (stakeholders), 81-109 (commitments), 110-134 (events), 815-922 (strata)
% src: erisml-compiler src/erisml_compiler/runtime/scene.py:48-66 (event types that drive machines; none is a standing event type)
% src: erisml-compiler src/erisml_compiler/runtime/strata.py:124-144 (no gate fires on consent_refused, role_assigned, authority_delegated or activity_started in this scene)

\section{One step moves strata, then machines, then norms}\label{sec:runtime-step}

The method \code{step} takes one event (Figure~\ref{fig:runtime-step}). An event given as a plain
dictionary is first built into an event record, with an identifier \code{live}$k$ and, when it has
no time index, the runtime's current time index plus one. The runtime's time index is the largest
seen so far. The event is appended to the run, and its position $k$ in the run is its index from
then on.
% src: erisml-compiler src/erisml_compiler/runtime/scene.py:192-199

\begin{figure}[t]
\centering
% Chapter 7. One step of the scene runtime: an event in, strata and machines step, a snapshot out.
% src: erisml-compiler src/erisml_compiler/runtime/scene.py:192-238 (step), 241-299 (in_force, snapshot)
% src: erisml-compiler src/erisml_compiler/runtime/strata.py:124-144 (cross)
% src: erisml-compiler src/erisml_compiler/runtime/agent.py:5-13, 299-309 (who supplies events)
\begin{tikzpicture}[house]
\def\tw{27mm}
% inputs
\node[panel=Grey, text width=24mm, minimum height=13mm] (cls) at (0,6mm) {\textcolor{hsGrey}{\bfseries classifier}\\{\scriptsize\color{hsTextMid}perceived types\\only}};
\node[panel=Red, text width=24mm, minimum height=13mm] (sys) at (0,-17mm) {\textcolor{hsRed}{\bfseries system}\\{\scriptsize\color{hsTextMid}rulings, replies,\\actions performed}};
\node[panel=Purple, text width=17mm, minimum height=13mm] (ev) at (29mm,-5.5mm) {\textcolor{hsPurple}{\bfseries event} $e_k$\\{\scriptsize\color{hsTextMid}type, content,\\actor, target}};
\draw[untrusted] (cls.east) -- (ev.west);
\draw[grant] (sys.east) -- (ev.west);
% row 1, left to right
\foreach \i/\c/\x/\name/\what/\det in {
  1/Teal/57/Strata/{each gate that\\fires on the event}/{first match wins\\records \texttt{entered\_at}\\a crossing if it moved},
  2/Green/90/Machines/{consent, legitimacy,\\commitment}/{each steps on its\\mapped event types},
  3/Orange/123/Defeasibility/{override or restore\\a commitment}/{a named condition\\overrides, only an\\oversight event restores}} {
  \node[panel=\c, text width=\tw, minimum height=38mm, anchor=north] (p\i) at (\x mm, 12mm) {};
  \node[step=\c] (s\i) at ([yshift=-5mm]p\i.north) {\i};
  \node[stepname=\c, below=0.8mm of s\i] (n\i) {\name};
  \node[desc, below=0.2mm of n\i, text width=\tw] (d\i) {\what};
  \node[detail, below=1mm of d\i, text width=\tw-2mm] {\det};
}
% row 2, right to left
\foreach \i/\c/\x/\name/\what/\det in {
  4/Amber/123/Discharge/{if the event is\\\texttt{action\_performed}}/{each obligation on\\that action records\\the event index},
  5/Blue/90/Norms in force/{commitment source\\and every token}/{prohibitions win ties\\against permissions}} {
  \node[panel=\c, text width=\tw, minimum height=38mm, anchor=north] (p\i) at (\x mm, -33mm) {};
  \node[step=\c] (s\i) at ([yshift=-5mm]p\i.north) {\i};
  \node[stepname=\c, below=0.8mm of s\i] (n\i) {\name};
  \node[desc, below=0.2mm of n\i, text width=\tw] (d\i) {\what};
  \node[detail, below=1mm of d\i, text width=\tw-2mm] {\det};
}
\node[panel=Grey, text width=40mm, minimum height=38mm, anchor=north] (snap) at (46mm, -33mm) {};
\node[stepname=Grey, below=2mm of snap.north] (sn) {Snapshot};
\node[detail, below=1mm of sn, text width=37mm, align=left] {\texttt{prohibited}\quad\texttt{allowed}\\\texttt{obliged}\quad\texttt{obliged\_tiers}\\\texttt{machines}\quad\texttt{reasons}\\\texttt{norms\_in\_force}\\\texttt{crossings}};
\draw[flow] (ev.east) -- (ev.east -| p1.west);
\draw[flow] (p1.east) -- (p1.east -| p2.west);
\draw[flow] (p2.east) -- (p2.east -| p3.west);
\draw[flow] (p3.south) -- (p4.north);
\draw[flow] (p4.west) -- (p4.west -| p5.east);
\draw[flow] (p5.west) -- (p5.west -| snap.east);
\end{tikzpicture}

\caption{One step of the scene runtime, from an event to a snapshot. An event reaches the runtime
from the classifier, which may propose only perceived types, or from the system, the only source
that can lead into an authority stratum.}
\label{fig:runtime-step}
\end{figure}

The step then runs five stages in a fixed order.
\begin{enumerate}
\item Every stratification is offered the event. The first of its gates that fires sets its state and
records the event's index as the moment the stratum was established. If the state changed, the
runtime records a crossing.
\item The consent, legitimacy and commitment machines step on the event types mapped to them.
\item Each commitment is checked for defeasibility. An active commitment whose defeasibility
condition now holds is overridden, and an overridden commitment is restored by an oversight event.
\item If the event is \code{action\_\allowbreak{}performed}, every obligation on the performed action records
the event's index as its discharge point.
\item The runtime computes the snapshot.
\end{enumerate}
The order matters. Strata step first, so a ruling that moves the situation into an emergency
overrides the privacy promise in the same step, because the promise's defeasibility condition reads
that stratum.
% src: erisml-compiler src/erisml_compiler/runtime/scene.py:201-238
% src: gtc-prototype twin/scene/margaret_home.erisml:100, 778 (privacy_promise defeasible on stratum:situation=emergency)

The snapshot carries the time index, the event, the state of every machine and stratification, the
identifiers of the norms in force, the prohibited, obliged and allowed actions, the norms behind each
of them, the crossings of this step, and the most urgent tier of each obliged action.
% src: erisml-compiler src/erisml_compiler/runtime/scene.py:70-103

\section{Three small machines hold commitments, consent and legitimacy}\label{sec:runtime-machines}

Each machine is a deterministic table from a state and an event tag to a next state
(Figure~\ref{fig:runtime-machines}). A tag the table does not list for the current state leaves the
state unchanged. Absorbing states ignore every tag. The commitment machine has seven states, and its
five terminal states are absorbing. The consent machine has four, and coerced and withdrawn are
absorbing. The legitimacy machine has six, and void is absorbing.
% src: erisml-compiler src/erisml_compiler/fsm/commitment_fsm.py:33-91
% src: erisml-compiler src/erisml_compiler/fsm/consent_fsm.py:13-37
% src: erisml-compiler src/erisml_compiler/fsm/legitimacy_fsm.py:22-74

\begin{figure}[t]
\centering
% Chapter 7. The three moral state machines, with the transitions their tables define.
% The red edge is the transition that lifts a promise's prohibitions. Dashed edges are in the
% machine's table but the scene runtime never sends their event tag. Double-bordered states are
% absorbing.
% src: erisml-compiler src/erisml_compiler/fsm/commitment_fsm.py:37-71
% src: erisml-compiler src/erisml_compiler/fsm/consent_fsm.py:18-21, 30-37
% src: erisml-compiler src/erisml_compiler/fsm/legitimacy_fsm.py:33-52, 67-74
% src: erisml-compiler src/erisml_compiler/runtime/scene.py:48-66, 214-230 (which tags the runtime sends)
\begin{tikzpicture}[house, lbl/.style={font=\sffamily\tiny, text=hsTextMid}, e/.style={->, draw=hsGreyLine, line width=0.9pt}]
% ---------------- commitment
\node[panel=Green, minimum width=158mm, minimum height=37mm, anchor=north west] (pc) at (0,0) {};
\node[stepname=Green, anchor=north west] at ([shift={(2mm,-1.5mm)}]pc.north west) {Commitment};
\node[stratum] (act) at (26mm,-14mm) {active};
\node[stratum, text width=20mm, align=center] (abd) at (76mm,-14mm) {active but\\defeasible};
\node[absorbing] (ful) at (132mm,-8mm) {fulfilled};
\node[absorbing] (vio) at (132mm,-19mm) {violated};
\node[absorbing] (exp) at (132mm,-30mm) {expired};
\node[absorbing] (voi) at (26mm,-30mm) {void};
\node[absorbing] (def) at (76mm,-30mm) {defeated};
\draw[grant] (act.north east) to[bend left=14] node[lbl, above]{condition triggered} (abd.north west);
\draw[e] (abd.west) -- node[lbl, below]{resolved in favour} (act.east);
\draw[e] (abd.east) -- (ful.west);
\draw[e] (abd.east) -- (vio.west);
\draw[e] (abd.east) -- (exp.west);
\draw[e, dashed] (act) -- node[lbl, left]{legitimacy collapses} (voi);
\draw[e, dashed] (abd) -- node[lbl, right]{overrides} (def);
\node[lbl, align=center] at (104mm,-35mm) {fulfilling, violating, expiration, from either live state};
% ---------------- consent
\node[panel=Purple, minimum width=76mm, minimum height=30mm, anchor=north west] (pk) at (0,-40mm) {};
\node[stepname=Purple, anchor=north west] at ([shift={(2mm,-1.5mm)}]pk.north west) {Consent};
\node[stratum, minimum width=14mm] (no) at (13mm,-57mm) {not obtained};
\node[stratum, minimum width=14mm] (ob) at (41mm,-50mm) {obtained};
\node[absorbing, minimum width=12mm] (co) at (41mm,-64mm) {coerced};
\node[absorbing, minimum width=12mm] (wd) at (65mm,-50mm) {withdrawn};
\draw[e] (no) -- node[lbl, above left]{given} (ob);
\draw[e] (no) -- node[lbl, below left]{coerced assent} (co);
\draw[e] (ob) -- node[lbl, above]{withdrawn} (wd);
\draw[e, dashed] (ob) -- node[lbl, right]{revisit} (co);
% ---------------- legitimacy
\node[panel=Blue, minimum width=79mm, minimum height=30mm, anchor=north west] (pl) at (79mm,-40mm) {};
\node[stepname=Blue, anchor=north west] at ([shift={(2mm,-1.5mm)}]pl.north west) {Legitimacy};
\node[stratum, minimum width=12mm] (fl) at (94mm,-57mm) {legitimate};
\node[stratum, minimum width=12mm] (df) at (114mm,-49mm) {defeasible};
\node[stratum, minimum width=12mm] (cr) at (132mm,-57mm) {coercive};
\node[stratum, minimum width=12mm] (fr) at (114mm,-65mm) {fraudulent};
\node[stratum, minimum width=12mm] (ty) at (149mm,-49mm) {tyrannical};
\node[absorbing, minimum width=10mm] (vd) at (149mm,-65mm) {void};
\draw[e] (fl) -- (df); \draw[e] (df) -- (cr); \draw[e] (fl) -- (cr);
\draw[e] (fl) -- (fr); \draw[e] (df) -- (fr);
\draw[e] (cr) -- (ty); \draw[e] (cr) -- (vd); \draw[e] (ty) -- (vd); \draw[e] (fr) -- (vd);
\draw[e] (df.west) to[bend right=25] node[lbl, above left]{restored} (fl.north);
\end{tikzpicture}

\caption{The three moral state machines as their tables define them. Dashed transitions exist in
the tables, but the scene runtime never sends their tags, so it never reaches defeated or void
through a commitment.}
\label{fig:runtime-machines}
\end{figure}

The runtime maps scene event types to tags. Three consent types, \code{consent\_given},
\code{consent\_\allowbreak{}withdrawn} and \code{coerced\_\allowbreak{}assent}, step the machine of the event's actor. Six
legitimacy types step the machine of the event's target. Three commitment types step the machine
whose identifier is the event's content. The consent tag for a coerced revisit and the commitment
tags for an override and for a legitimacy collapse are in the tables but not in the runtime's maps.
The legitimacy table also differs from its docstring. The table lets a fully legitimate or
defeasible authority become fraudulent and lets a coercive one become void, which the docstring
does not list. This chapter follows the tables.
% src: erisml-compiler src/erisml_compiler/runtime/scene.py:48-66, 207-212
% src: erisml-compiler src/erisml_compiler/fsm/consent_fsm.py:1-21
% src: erisml-compiler src/erisml_compiler/fsm/commitment_fsm.py:44-71
% src: erisml-compiler src/erisml_compiler/fsm/legitimacy_fsm.py:1-52

Margaret's scene declares none of these twelve event types in its vocabulary. In the twin the
commitment machines move only through defeasibility and oversight, and the consent and legitimacy
machines keep their initial states.
% src: grep of the twelve type names over gtc-prototype twin/scene/margaret_home.erisml and twin/*.py returned nothing (2026-10-03)

\section{A condition overrides a commitment and only oversight restores it}\label{sec:runtime-defeasible}

After the machines step, the runtime checks every commitment. An active commitment is overridden
when any of its named defeasibility conditions holds, read from the point of its last restoration.
An overridden commitment is restored when the event's type is in the scene's \code{oversight} list
and the event names the commitment by content or target, or names nothing. Restoration records the
event's index, and only later events can override the commitment again.
% src: erisml-compiler src/erisml_compiler/runtime/scene.py:214-230
% src: erisml-compiler tests/test_scene_runtime.py:98-114

In Margaret's home the privacy promise is overridden when the situation stratification enters
emergency, and its three prohibitions are lifted while it is overridden. Two of the four oversight
types can restore it. The centre's confirmation that privacy is restored and the lapse of evidence
declare no content, so an event of either type that names nothing restores any overridden
commitment. The clearing of a visitor is declared with the single content \code{visitor}, and the
clearing of devices with \code{devices}. Neither names a commitment, so neither restores one. The
scene's comment says the same of the first.
% src: gtc-prototype twin/scene/margaret_home.erisml:100, 778, 802-804, 1120, 1133, 1139, 1151
% src: erisml-compiler src/erisml_compiler/runtime/scene.py:224-230

The freshness rule of the next section applies here too. After a restoration, the emergency
stratum counts toward a new override only if a gate establishes it again after the restoration. A
ruling of elevate that arrives after the restoration overrides the promise again.
% src: erisml-compiler src/erisml_compiler/runtime/scene.py:180-188, 217-223

\section{Prohibitions win}\label{sec:runtime-norms}

Let $N$ be the scene's norms and $A$ the actions in its capability list. For a norm $n$ write
$\pi(n)$ for its priority tier and $d_n$ for its discharge point, which is $0$ until an action
discharges it. A norm is lifted when it is defeasible, its source is \code{commitment:}$c$ for a
commitment $c$ of the scene, and that commitment's machine is not active. The set of norms in force
after event $k$ is
\begin{equation}
F_k = \bigl\{\, n \in N \bigm| n \text{ is not lifted} \;\wedge\; \forall t \in T(n)\;
\mathrm{holds}_{s(n)}(t) \,\bigr\},
\label{eq:in-force}
\end{equation}
where $T(n)$ is the norm's condition list, and $s(n) = d_n$ for an obligation and $s(n) = 0$ for any
other norm.
% src: erisml-compiler src/erisml_compiler/runtime/scene.py:241-251

A norm applies to an action $a$ when the scene has no agent or the norm's actor is the agent, and the
norm's action is $a$ or is \code{elevated\_\allowbreak{}action} with $a$ an elevated capability. For each
$a \in A$ let $P_a$ be the prohibitions in $F_k$ that apply to $a$, and let $R_a$ be the permissions
and exceptions in $F_k$ that apply to $a$ by name. Then $a$ is prohibited when
\begin{equation}
P_a \neq \emptyset \;\wedge\; \Bigl( R_a = \emptyset \;\vee\;
\min_{n \in P_a} \pi(n) \;\le\; \min_{n \in R_a} \pi(n) \Bigr).
\label{eq:prohibited}
\end{equation}
A permission lifts a prohibition only from a strictly more urgent tier. On equal tiers the
prohibition wins. A permission written against \code{elevated\_\allowbreak{}action} lifts nothing. The allowed
actions are the capabilities that are not prohibited.
% src: erisml-compiler src/erisml_compiler/runtime/scene.py:253-288
% src: erisml-compiler tests/test_scene_runtime.py:128-141

The obliged actions are the actions of the obligations in $F_k$ that bind the agent. The obliged set
is not intersected with the capabilities and is not filtered by the prohibitions. An action can be
obliged and prohibited in the same snapshot. The scene agent's chooser removes that conflict in the
direction of the prohibition, because it offers the model only obligations that are allowed and
falls back only to an allowed obligation (Section~\ref{sec:runtime-agent}). Prohibitions therefore
win twice, once against permissions in Equation~\eqref{eq:prohibited} and once against obligations
in the chooser. For each prohibited or obliged action the snapshot records the identifiers of the
norms behind it.
% src: erisml-compiler src/erisml_compiler/runtime/scene.py:282-298
% src: erisml-compiler src/erisml_compiler/runtime/agent.py:243-244, 269-275

Margaret's emergency call shows the rule at work (Listing~\ref{lst:scene-ems}). Outside an
emergency and without the governor's authorization the prohibition \code{n6p} is in force at tier 0
and no permission exists, so the call is prohibited. Under a ruling of elevate the situation
stratification enters emergency. The prohibition leaves $F_k$ and the obligation \code{n6} enters it.
% src: gtc-prototype twin/scene/margaret_home.erisml:187-208, 778, 785, 873-880

\paragraph{Urgency reaches the chooser.}
Each obliged action carries the most urgent tier among the obligations in force that oblige it. The
snapshot stores these tiers, and the method \code{obliged\_\allowbreak{}by\_\allowbreak{}priority} orders the obliged actions
by tier and then by name. The chooser receives the allowed obligations in that order, each with its
tier, and is told that tier 0 is the most urgent. Its fallback takes the first. Before this change
the model saw obligations in alphabetical order with no tier, and in a twin development run a
referral of fraud to the monitoring centre lost to check-ins. The obliged set itself did not change.
% src: erisml-compiler src/erisml_compiler/runtime/scene.py:84-89, 282-298
% src: erisml-compiler src/erisml_compiler/runtime/agent.py:42-50, 243-251, 270-274
% src: erisml-compiler PR #38 (merged), "Found in the gtc twin (dev12b d24)"
% src: erisml-compiler tests/test_scene_agent.py:248 (test_the_chooser_sees_obligations_most_urgent_first_and_falls_back_to_the_most_urgent)

\section{An action discharges an obligation until its stratum is established again}\label{sec:runtime-discharge}

An obligation is discharged by doing what it obliges. When the event \code{action\_\allowbreak{}performed} with
content $a$ is stepped at index $k$, every obligation whose action is $a$ records $d_n = k$, whether
or not it is in force at that moment. From then on its tokens are read from $d_n$
(Equation~\ref{eq:in-force}). An \code{event:} token over the trigger counts only a new event after
the action, so a fresh inactivity report obliges a fresh check-in. An obligation with no conditions
cannot be discharged and is in force at every step.
% src: erisml-compiler src/erisml_compiler/runtime/scene.py:232-235, 250-251
% src: erisml-compiler tests/test_scene_runtime.py:117-125

A stratum is a state, not an event, and this broke discharge. A token that read only the current
stratum ignored the discharge point, so an obligation that rested on a stratum stayed in force for
as long as the stratum lasted. The twin found this when the corroborated emergency moved from a
\code{latest:} token over the governor's ruling to the situation stratification. The robot remained
obliged to call emergency services after it had called them. Pull request 36 of the compiler
changed the rule.
% src: erisml-compiler PR #36 (merged) description
% src: erisml-compiler commit 4f9a224

Let $\varepsilon_\sigma$ be the index of the event at which a gate of stratification $\sigma$ last
fired, and let $\sigma$ be in state $x_\sigma$. A \code{stratum:}$\sigma$\code{=}$q$ token evaluated
from starting point $s$ holds when
\begin{equation}
x_\sigma = q \;\wedge\; \bigl(\, \nu \;\vee\; s = 0 \;\vee\; \varepsilon_\sigma > s \,\bigr),
\label{eq:stratum-fresh}
\end{equation}
where $\nu$ is true when the token sits under an odd number of \code{not:} prefixes. A positive
token counts after a discharge only if a gate established the stratum after it. Under a negation the
token reads the stratum as it stands now, so \code{not:} keeps the meaning ``not in that stratum
now''. The polarity passes through named conditions, so a \code{not:cond:} over a condition defined
by a stratum token also reads the stratum now. Margaret's check-in obligation after a fall carries
\code{not:cond:corroborated\_\allowbreak{}life\_\allowbreak{}threatening\_\allowbreak{}emergency} and so stops as soon as the situation is an
emergency, whatever was discharged before.
% src: erisml-compiler src/erisml_compiler/runtime/scene.py:151-162, 173-176, 180-188
% src: gtc-prototype twin/scene/margaret_home.erisml:248-256 (n4c), 778

A gate that fires records the index even when it leaves the stratum where it was. Fresh evidence for
the same stratum therefore owes the obligation again without any crossing. Three tests fix the rule,
namely discharge until the stratum is entered again, a negated stratum read now after a discharge,
and fresh evidence for the same stratum.
% src: erisml-compiler src/erisml_compiler/runtime/strata.py:124-144
% src: erisml-compiler tests/test_scene_strata.py:176-225

Figure~\ref{fig:runtime-discharge} works the rule by hand for the obligation \code{p9v}, which obliges
the robot to speak, to warn Margaret, while an animal at hand is venomous.
% src: gtc-prototype twin/scene/margaret_home.erisml:570-578 (p9v, stratum:animal_standing=venomous)

\begin{figure}[t]
\centering
% Chapter 7. Obligation discharge with a stratum token, worked by hand from the code (not a recorded run).
% Norm p9v of Margaret's home: obligation speak, condition stratum:animal_standing=venomous.
% Gates a_venomous (animal_seen=venomous, from none/visiting/pest/wild) and a_clear (animal_clear).
% src: gtc-prototype twin/scene/margaret_home.erisml:570-578 (p9v), 855-865 (animal_standing)
% src: erisml-compiler src/erisml_compiler/runtime/scene.py:180-188 (fresh rule), 232-235 (discharge)
% src: erisml-compiler src/erisml_compiler/runtime/strata.py:124-144 (entered_at set when a gate fires)
% src: erisml-compiler PR #36 (before: a stratum token ignored the discharge point)
\begin{tikzpicture}[house, row/.style={font=\sffamily\scriptsize\bfseries, text=hsTextLight, anchor=east, align=right},
  cell/.style={font=\sffamily\scriptsize, text=hsText, align=center}]
\def\cw{27mm}
\foreach \j/\ev in {1/{\texttt{animal\_seen}\\\texttt{=venomous}}, 2/{\texttt{action\_}\\\texttt{performed}\\\texttt{=speak}}, 3/{\texttt{activity\_}\\\texttt{started}}, 4/{\texttt{animal\_clear}}, 5/{\texttt{animal\_seen}\\\texttt{=venomous}}} {
  \node[panel=Purple, text width=\cw-8mm, minimum height=11mm, inner sep=2pt, font=\sffamily\scriptsize] (h\j) at (\j*\cw, 0) {\ev};
  \node[step=Purple, anchor=south] at ([yshift=-1mm]h\j.north) {\j};
}
\node[row] at (0.45*\cw, 0) {event};
% stratum row
\foreach \j/\st/\sty in {1/venomous/stratum, 2/venomous/stratum, 3/venomous/stratum, 4/none/stratum, 5/venomous/stratum} {
  \node[\sty, minimum width=18mm] (a\j) at (\j*\cw, -12mm) {\st};
}
\node[row] at (0.45*\cw, -12mm) {\texttt{animal\_standing}};
% entered_at and discharge point
\foreach \j/\en/\d in {1/1/0, 2/1/2, 3/1/2, 4/4/2, 5/5/2} {
  \node[cell] at (\j*\cw, -20mm) {\en};
  \node[cell] at (\j*\cw, -26mm) {\d};
}
\node[row] at (0.45*\cw, -20mm) {\texttt{entered\_at}};
\node[row] at (0.45*\cw, -26mm) {discharge point};
% owed rows
\foreach \j/\now/\old in {1/Green/Green, 2/Grey/Green, 3/Grey/Green, 4/Grey/Grey, 5/Green/Green} {
  \node[panel=\now, minimum width=21mm, minimum height=6mm, inner sep=1pt, font=\sffamily\scriptsize] at (\j*\cw, -35mm) {};
  \node[panel=\old, minimum width=21mm, minimum height=6mm, inner sep=1pt, font=\sffamily\scriptsize, dashed] at (\j*\cw, -44mm) {};
}
\foreach \j/\a/\b in {1/owed/owed, 2/discharged/still owed, 3/discharged/still owed, 4/not in stratum/not in stratum, 5/owed/owed} {
  \node[font=\sffamily\scriptsize, text=hsText] at (\j*\cw, -35mm) {\a};
  \node[font=\sffamily\scriptsize, text=hsTextMid] at (\j*\cw, -44mm) {\b};
}
\node[row] at (0.45*\cw, -35mm) {\texttt{speak} obliged\\since PR \#36};
\node[row] at (0.45*\cw, -44mm) {before PR \#36};
\foreach \j/\k in {1/2, 2/3, 3/4, 4/5} { \draw[flow] (h\j.east) -- (h\k.west); }
\end{tikzpicture}

\caption{The obligation to warn Margaret about a venomous animal through five events, worked from
the code and not from a recorded run. Before pull request 36 the warning stayed owed after it was
given, for as long as the animal stayed venomous.}
\label{fig:runtime-discharge}
\end{figure}

A \code{state:} token over a stratum, written \code{state:stratum:}$\sigma$\code{=}$q$, reads the
same state but bypasses the freshness rule, because the runtime's \code{state:} branch ignores the
starting point. Margaret's scene uses no \code{state:} token.
% src: erisml-compiler src/erisml_compiler/runtime/scene.py:138-143, 177-179
% src: gtc-prototype twin/scene/margaret_home.erisml (no "state:" token; grep 2026-10-03)

\section{Strata are stepped by gates and loaded under two checks}\label{sec:runtime-strata}

A gate has an identifier, a trigger type, an optional trigger content, a target state, a tuple of
source states and a boundary type. It fires on an event when the event's type is the trigger type,
the event's content equals the trigger content if one is given, and the stratification is in one of
the source states, or in any state when none is listed. A stratification offers each event to its
gates in declared order, and the first that fires wins. Gates are edge-triggered. Only the event
being stepped can fire one, and history never fires it again. A crossing record names the
stratification, the gate, the boundary type and the source and target states.
% src: erisml-compiler src/erisml_compiler/runtime/strata.py:17-25, 46-60, 124-144
% src: erisml-compiler tests/test_scene_strata.py:72-95

The boundary types follow Geometric Ethics \cite{bond2026geometricethics}. A \code{threshold} gate is
Type~I, a measured quantity crossing a value. A \code{phase} gate is Type~II, a change of phase. An
absorbing state is Type~III. Type~IV constraint surfaces are not states. They are the scene's non-defeasible
prohibitions.
% src: erisml-compiler src/erisml_compiler/runtime/strata.py:17-22, 43

The loader refuses a scene that breaks any of its rules, so such a scene never runs. A stratification
must declare states, and its initial state must be one of them. Its authority and absorbing states
must be declared, and its initial state may not be an authority state. Every gate must have a
trigger, name only declared states, use one of the two boundary types and fire on a declared event
type. Two further checks carry the containment argument.
% src: erisml-compiler src/erisml_compiler/runtime/strata.py:26-33, 63-122
% src: erisml-compiler tests/test_scene_strata.py:117-165
\begin{description}
\item[Authority.] A gate whose target is an authority state must fire on an event type declared
\code{source: system}. A model's reading of the world cannot lead into a state that grants
something.
\item[Absorption.] A gate that can leave an absorbing state must fire on a human-oversight event. The
states a gate can leave are its source states, or every state when it lists none, less its target.
\end{description}
% src: erisml-compiler src/erisml_compiler/runtime/strata.py:26-33, 113-122

In Margaret's home both authority states of the situation are entered only on the governor's ruling.
The reporting duty has no authority state, so a perceived suspicion of a crime may enter it. Adding a
duty grants nothing. The visitor's hostile state is left only on \code{visitor\_\allowbreak{}cleared}, and a
compromised device only on \code{device\_\allowbreak{}cleared}, both oversight types.
% src: gtc-prototype twin/scene/margaret_home.erisml:804, 820-831, 873-880, 893-907, 916-922

\paragraph{The trigger key.}
A gate's trigger was first written under the key \code{on}. YAML 1.1 reads a bare \code{on} as the
Boolean true, so a scene written as \verb|{on: event}| produced a gate keyed \code{True} with an empty
trigger. The loader then refused it as an undeclared event type with an empty name, which hid the
cause, and the twin's scene hit exactly this. Pull request 29 renamed the key to \code{trigger}, the
book's triggering feature, and the loader now refuses a gate without one with a message that names
the YAML trap. A test parses \code{on:} with PyYAML and checks the refusal.
% src: erisml-compiler PR #29 (merged) description
% src: erisml-compiler src/erisml_compiler/runtime/strata.py:17-19, 87-90
% src: erisml-compiler tests/test_scene_strata.py:167-174

\section{The agent's two models are isolated from each other}\label{sec:runtime-agent}

The runtime decides nothing about the world. The scene agent connects it to two models
(Figure~\ref{fig:runtime-agent}). The observation classifier turns perception facts into events of
the scene's vocabulary. The action chooser picks one action from the allowed set. Neither model can
change the norms, the machines or the allowed set. Both receive the scene's role, so one runtime can
serve every agent of a scene with several parties.
% src: erisml-compiler src/erisml_compiler/runtime/agent.py:1-22, 56-57

\begin{figure}[t]
\centering
% Chapter 7. The isolated scene agent: two models, the checks between them, and the runtime they read.
% Dashed grey edges carry what a model wrote. The red edge carries system events, the only input
% that can move an authority stratum.
% src: erisml-compiler src/erisml_compiler/runtime/agent.py:84-201 (ObservationClassifier)
% src: erisml-compiler src/erisml_compiler/runtime/agent.py:204-275 (ActionChooser, canonical view, fallback)
% src: erisml-compiler src/erisml_compiler/runtime/agent.py:278-309 (SceneAgent.decide, record)
\begin{tikzpicture}[house]
\def\tw{23mm}
\foreach \i/\c/\x/\name/\what/\det in {
  1/Grey/0/Facts/{perception, any JSON}/{readings, poses,\\words heard,\\text on a screen},
  2/Orange/35.5/Classifier/{a model proposes\\events}/{scene text,\\vocabulary,\\last 8 events,\\facts},
  3/Blue/71/Checks/{reject, snap,\\quarantine}/{bad types and\\contents rejected,\\actors snapped,\\free text\\quarantined},
  4/Green/106.5/Runtime/{steps the canonical\\events}/{allowed,\\obliged by urgency,\\prohibited,\\machines},
  5/Orange/142/Chooser/{a model proposes\\one action}/{last 16 events,\\moral state,\\never the facts}} {
  \node[panel=\c, text width=\tw, minimum height=42mm, anchor=north] (p\i) at (\x mm, 0) {};
  \node[step=\c] (s\i) at ([yshift=-5mm]p\i.north) {\i};
  \node[stepname=\c, below=0.8mm of s\i] (n\i) {\name};
  \node[desc, below=0.2mm of n\i, text width=\tw] (d\i) {\what};
  \node[detail, below=1mm of d\i, text width=\tw-2mm] {\det};
}
\draw[flow] (p1.east) -- (p2.west);
\draw[untrusted] (p2.east) -- node[flowlabel, above]{events} (p3.west);
\draw[flow] (p3.east) -- (p4.west);
\draw[flow] (p4.east) -- (p5.west);
\node[panel=Red, text width=32mm, below=9mm of p4, xshift=-8mm] (sys) {\textcolor{hsRed}{\bfseries system events}\\{\scriptsize\color{hsTextMid}rulings, channel replies,\\actions performed}};
\draw[grant] (sys.north -| p4.south) -- (p4.south);
\node[panel=Teal, text width=31mm, below=9mm of p5, xshift=1mm] (gate) {\textcolor{hsTeal}{\bfseries allowed?}\\{\scriptsize\color{hsTextMid}else the most urgent\\allowed obligation,\\else the default}};
\draw[untrusted] (p5.south) -- node[flowlabel, right]{proposal} (gate.north -| p5.south);
\node[panel=Grey, text width=30mm, below=9mm of p2] (rej) {\textcolor{hsGrey}{\bfseries reports}\\{\scriptsize\color{hsTextMid}rejected events, snapped\\actors, quarantined text}};
\draw[flow] (p3.south) |- (rej.east);
\end{tikzpicture}

\caption{The isolated scene agent. What the two models write passes through checks before it
reaches the runtime or the robot, and the chooser never sees the perception facts.}
\label{fig:runtime-agent}
\end{figure}

\paragraph{The classifier.}
The classifier's vocabulary is the scene's event types less the system types. It receives the
scene's plain-language description, that vocabulary, the last eight events and the facts. Each
event it proposes is checked. An event is rejected and reported, never stepped, when it is not an
object with a string type, when its actor, target or content is not a string, when its conditions
are not a list of strings, when its type is a system type, when its type is undeclared, or when its
content is outside the type's declared list. A comment in the code records why the field types are
checked. A model once returned a list as the content of a free-content type, and the whole decision
failed.
% src: erisml-compiler src/erisml_compiler/runtime/agent.py:105-140, 153-172

Isolation adds a second layer, which the code calls the boundary of mandatory canonicalization. When
the scene declares actors in \code{extra.actors}, an event's actor must be one of them. Any other
actor is snapped by the canonicalizer to a declared one, or to \code{unknown}, and the snap is
recorded. An event type that
lists its actors keeps only those, and any other actor becomes the first listed. The target is
canonicalized like the actor. The event's conditions are dropped. Content for a type with no
declared content list is free text, and it is removed and recorded in a quarantine list, never
stepped. What is stepped is then canonical, with declared types, declared contents and declared
actors.
% src: erisml-compiler src/erisml_compiler/runtime/agent.py:84-95, 142-151, 183-200

The actor constraint came from the twin. In one development run, 30 sensor readings were classified
with the actor \code{unknown\_\allowbreak{}person} and 71 with \code{unknown}. Downstream they read as a stream of
strangers, and the chooser called the monitoring centre in scenarios where nobody was there. Pull
request 37 let an event type name its only actors, and Margaret's scene gives sensor readings the
actor \code{device} and media content the actor \code{media} (Listing~\ref{lst:scene-actor-types}).
% src: erisml-compiler PR #37 (merged) description, "dev11b d23, d25"
% src: erisml-compiler src/erisml_compiler/runtime/agent.py:186-194
% src: erisml-compiler tests/test_scene_agent.py:167
% src: gtc-prototype twin/scene/margaret_home.erisml:1121-1122

\paragraph{The chooser.}
The chooser has two views. In the facts view it receives the perception facts. In the canonical view
it never sees them. It receives the last sixteen stepped events reduced to type, actor, target and
content, and from the caller's context only two keys, \code{governor\_\allowbreak{}ruling} and
\code{requested}. Raw text
in the world, such as a television, a message or a stranger's words, then cannot reach the model
that picks the action. In both views it also receives the scene's description, the allowed
obligations with their tiers, the default action, the allowed capabilities with their descriptions,
the prohibited actions and the state of every machine.
% src: erisml-compiler src/erisml_compiler/runtime/agent.py:204-258

A choice outside the allowed set is rejected and recorded. The chooser then falls back to the most
urgent allowed obligation, else to the default action if it is allowed, else to the first allowed
action. The decision records the rejected proposal and marks the fallback.
% src: erisml-compiler src/erisml_compiler/runtime/agent.py:259-275
% src: erisml-compiler tests/test_scene_agent.py:49-67

\paragraph{One decision.}
The method \code{decide} classifies the facts, steps each accepted event in turn and asks the chooser
for an action on the last snapshot. Events that do not come from perception, such as a ruling or an
action performed, are stepped through \code{record} and never pass through the classifier. In the
twin the robot's agent is isolated and uses the compiler's registry canonicalizer. The monitoring
centre and the dispatcher are agents of their own scenes and run without isolation.
% src: erisml-compiler src/erisml_compiler/runtime/agent.py:278-309
% src: gtc-prototype twin/brain.py:185-194 (Desk builds SceneAgent without isolation), 304-305 (centre and EMS desks), 312-320 (robot, isolated=True, RegistryCanonicalizer)

These patterns have names in prior work. A model that reads untrusted text and returns only
constrained output is the quarantined model of the dual-model pattern \cite{willison2023dualllm}.
Tracking provenance outside the model so that it cannot grant itself capabilities is the design
principle of CaMeL \cite{debenedetti2026camel}. A model that may only select among permitted actions
is the action-selector pattern \cite{beurerkellner2025patterns}. A verified component that removes
unsafe choices of a learner is a shield \cite{bloem2015shield,alshiekh2018shielding}. The canonical
view follows the mandatory canonicalization of the No Escape argument \cite{bond2026noescape}.
% src: gtc-prototype paper/aies2027/prior_art.md:86-91, 126-139
% src: erisml-compiler src/erisml_compiler/runtime/agent.py:208-216 (cites erisml-lib docs/papers/foundations/no_escape.tex)

\section{A process is a checked view, exported to BPMN and imported under the same checks}\label{sec:runtime-processes}

The process loader turns each entry of \code{extra.processes} into nodes and flows and refuses any
process that fails one of its checks (Figure~\ref{fig:runtime-process}). Every node needs a unique
identifier that is a valid XML name, one of the six kinds and a declared lane. A task in the agent's
lane must name one of the scene's capabilities. A task in another lane must be marked external,
another party's work. A start or message event must be triggered by a declared event type, and a
message, being an authenticated channel, by a system type. A timer must name an interval the scene
declares, as a dotted path into \code{extra}, in seconds, minutes or hours. A flow's condition must be
evaluable in the scene. No flow may enter a start event or leave an end event. Every node must be
reachable from a start event, and every node but an end event must have a way on.
% src: erisml-compiler src/erisml_compiler/process/model.py:20-34, 139-251

\begin{figure}[t]
\centering
% Chapter 7. A scene's process: declared, checked, exported to BPMN 2.0, edited, imported under the same
% checks, and traced against a run.
% src: erisml-compiler src/erisml_compiler/process/model.py:1-34, 139-251
% src: erisml-compiler src/erisml_compiler/process/bpmn.py:1-18, 163-320
% src: erisml-compiler src/erisml_compiler/process/importer.py:1-25, 36-52, 90-99, 255-294
% src: erisml-compiler src/erisml_compiler/process/trace.py:1-10, 28-75
\begin{tikzpicture}[house]
\def\tw{26mm}
\foreach \i/\c/\x/\name/\what/\det in {
  1/Purple/0/Declared/{\texttt{extra.processes}}/{lanes, nodes,\\flows with\\condition tokens},
  2/Blue/33/Checked/{\texttt{model.load}}/{capabilities, types,\\system messages,\\timers from \texttt{extra},\\containment,\\reachability},
  3/Green/66/Exported/{\texttt{to\_bpmn}}/{BPMN 2.0 with\\diagram interchange\\\texttt{erisml:binding}\\token conditions},
  4/Grey/99/Edited/{any BPMN modeler}/{bpmn-js,\\Camunda Modeler},
  5/Orange/132/Imported/{\texttt{import\_bpmn}}/{refuses parallel and\\other gateways,\\scripts, foreign\\expressions, DTDs,\\edited durations}} {
  \node[panel=\c, text width=\tw, minimum height=41mm, anchor=north] (p\i) at (\x mm, 0) {};
  \node[step=\c] (s\i) at ([yshift=-5mm]p\i.north) {\i};
  \node[stepname=\c, below=0.8mm of s\i] (n\i) {\name};
  \node[desc, below=0.2mm of n\i, text width=\tw] (d\i) {\what};
  \node[detail, below=1mm of d\i, text width=\tw-2mm] {\det};
}
\foreach \a/\b in {1/2, 2/3, 3/4, 4/5} { \draw[flow] (p\a.east) -- (p\b.west); }
\draw[flow] (p5.north) .. controls +(0,12mm) and +(0,12mm) .. node[flowlabel, above]{the same checks again} (p2.north);
\node[panel=Teal, text width=60mm, below=8mm of p2, xshift=16mm] (tr) {\textcolor{hsTeal}{\bfseries Trace}\\{\scriptsize\color{hsTextMid}records of a run activate starts, messages and tasks\\gateways, timers and other parties' tasks fill the path\\a view of what happened, it decides nothing}};
\draw[flow] (p2.south) -- (p2.south |- tr.north);
\node[panel=Grey, text width=28mm, below=8mm of p5] (run) {\textcolor{hsGrey}{\bfseries agent records}\\{\scriptsize\color{hsTextMid}events stepped,\\actions performed}};
\draw[flow] (run.west) -- (run.west -| tr.east);
\end{tikzpicture}

\caption{A scene's process from declaration to diagram and back. An edited diagram re-enters the
scene only through the checks a declared process passes.}
\label{fig:runtime-process}
\end{figure}

The containment check applies to tasks whose capability is elevated or governed. A flow into such a
task must rest only on structural condition tokens, those that only system events can make true.
A flow with no condition that leaves a start or message event triggered by a system event is also
structural, because the authenticated trigger is the evidence. Margaret's escalation ladder uses both
forms. Its start event for a corroborated emergency is triggered by the governor's ruling of elevate
and flows straight into the call to emergency services. Its gateway after a ruling reaches the same
call on \code{latest:governor\_\allowbreak{}ruling=authorize\_\allowbreak{}ems}.
% src: erisml-compiler src/erisml_compiler/process/model.py:30-33, 210-219
% src: erisml-compiler tests/test_process_bpmn.py:196 (a governed task may follow a system-triggered event directly)
% src: gtc-prototype twin/scene/margaret_home.erisml:934, 952, 958, 977, 1162

\paragraph{Export.}
The function \code{to\_bpmn} writes a checked process as BPMN 2.0 XML \cite{omg2011bpmn} with diagram
interchange, so that the file opens in BPMN tools such as bpmn-js and Camunda Modeler. The tests
validate the output against the OMG's schemas. The process becomes one pool with a lane per party. A
task carries its capability in an \code{erisml:binding} extension element. A start event triggered by
a system event becomes a message start, and one triggered by a perceived event becomes a
conditional start whose condition is that event's token. A message becomes an intermediate message
catch, and a timer an intermediate timer catch with an ISO 8601 duration resolved from the scene. The
conditions on a gateway's flows are written as formal expressions in the language
\code{https://erisml.org/condition-token}. The layout is layered, with a node's column given by its
longest distance from a start and its row by its lane. A descriptor file lets modelers keep the
binding when they save.
% src: erisml-compiler src/erisml_compiler/process/bpmn.py:1-18, 28-36, 163-320
% src: erisml-compiler src/erisml_compiler/process/__init__.py:5-8 (erisml.json moddle descriptor)
% src: erisml-compiler tests/test_process_bpmn.py:1-18, 110

\paragraph{Import.}
A BPMN document is a proposal, and the scene decides. The importer reads the meaning back from the
places the exporter writes it and refuses what the runtime cannot honour. It refuses parallel,
inclusive, complex and event-based gateways, sub-processes, transactions, call activities, boundary
events, intermediate throw events, script tasks, business-rule tasks and receive tasks. It refuses a
condition in any expression language other than condition tokens. It refuses a timer that does not
name the scene's interval, and a timer whose diagram duration disagrees with that interval, so a
duration edited in a modeler is refused instead of silently ignored. It refuses a document that
declares a DTD or an entity, and a document over five million bytes. It then runs the same checks as
for a declared process.
% src: erisml-compiler src/erisml_compiler/process/importer.py:1-25, 36-52, 90-99, 226-232, 255-264, 275-294
% src: erisml-compiler tests/test_process_import.py:24-170

\paragraph{Trace.}
The trace maps an agent's record stream onto a process for a live view. An event activates a start or
message node whose trigger matches its type and content. An \code{action\_\allowbreak{}performed} record activates
the task for that action. Between two activated nodes the trace recovers the shortest path through
the nodes a run does not record, namely gateways, timers and other parties' tasks, up to a depth of
four. The trace decides nothing. A task appears in it only if the agent performed the action, and the
agent performed it only if the scene allowed it.
% src: erisml-compiler src/erisml_compiler/process/trace.py:1-10, 28-75

The use of BPMN to describe robot missions has precedent \cite{corradini2023bpmn}. Here it describes
an escalation of human authority and stays a view of the norms.
% src: gtc-prototype paper/aies2027/prior_art.md:170-175

\section{What the runtime does not do}\label{sec:runtime-limits}

The runtime keeps no clock. Time is the order of events, and a timer in a process names an interval
that code outside the runtime must measure and report as an event. The runtime does not check that an
event's type is declared. Checks happen where events enter and where gates and processes load. It
does not check that an obliged action is one of the agent's capabilities. It does not resolve a
conflict between an obligation and a prohibition. It reports both, and the chooser resolves it toward
the prohibition. It does not decide whether an allowed action is wise or safe for Margaret. In the
twin that judgement belongs to the governor and the \deme{} output gate of Chapters~6 and~9.
% src: erisml-compiler src/erisml_compiler/runtime/scene.py:192-238, 282-298
% src: erisml-compiler src/erisml_compiler/runtime/agent.py:16-18 (authority for elevated actions stays with the caller's gate)
% src: gtc-prototype twin/output_gate.py:1-5
